\ExerciseSection

\begin{enumerate}
\def\labelenumi{\arabic{enumi})}
\tightlist
\item
  设 \(f\) 是 \(\mathbf{S}_2\) 到 \(\mathbb{R}^4\) 中的映射，满足
\end{enumerate}

\[
f \left(x _ {1}, x _ {2}, x _ {3}\right) = \left(x _ {1} ^ {2} - x _ {2} ^ {2}, x _ {1} x _ {2}, x _ {1} x _ {3}, x _ {2} x _ {3}\right).
\]

这个函数在 \(S_{2}\) 的对径点处取相同的值。证明：过渡到商后，它定义 \(P_{2}\) 到 \(R^{4}\) 的一个子空间上的一个同胚。

同样证明：由下式定义的映射 \(g\) （ \(\mathbf{S}_3\) 到 \(\mathbf{R}^6\) 中的）

\[
g \left(x _ {1}, x _ {2}, x _ {3}, x _ {4}\right) = \left(x _ {1} ^ {2} - x _ {2} ^ {2}, x _ {1} x _ {2}, x _ {1} x _ {3} + x _ {2} x _ {4}, x _ {3} ^ {2} - x _ {4} ^ {2}, x _ {3} x _ {4}, x _ {1} x _ {4} - x _ {2} x _ {3}\right)
\]

在过渡到商后定义 \(P_{3}\) 到 \(R^{6}\) 中的一个同胚。

\begin{enumerate}
\def\labelenumi{\arabic{enumi})}
\setcounter{enumi}{1}
\item
  把射影空间 \(P_{n}\) 与第 1 小节命题 3 中所定义的 \(B_{n}\) 的商空间等同起来，并设 \(\varphi\) 表示 \(B_{n}\) 到 \(P_{n}\) 上的典范映射。拓扑空间 E 到 \(P_{n}\) 中的连续映射 f 称为非本质的，如果存在 E 到 \(B_{n}\) 中的连续映射 g，使得 \(f = \varphi \circ g\) ，反之的情形则称为本质的（参见 § 2 的习题 6）。证明：存在 \(P_{n}\) 的一致结构的一个围域 U，使得若 f 是拓扑空间 E 到 \(P_{n}\) 中的一个非本质映射，则每个连续映射 \(f'\) （ E 到 \(P_{n}\) 中的，且使 \((f(x), f'(x)) \in \mathbf{U}\) 对一切 \(x \in E\) 成立）也是非本质的。
\item
  若 \(\mathbf{S}_1\) 到 \(\mathbf{P}_n\) 中的一个连续映射是本质的（习题 2），则它不能延拓为 \(\mathbf{B}_2\) 到 \(\mathbf{P}_n\) 中的连续映射（参见 § 2 的习题 8）。
\item
  若 \(n > \mathbf{1}\) ，则存在本质映射 \(f\) （ \(\mathbf{S}_1\) 到 \(\mathbf{P}_n\) 中的），使得 \(f(\mathbf{S}_1) \neq \mathbf{P}_n\) {[}取 \(f(\mathbf{S}_1)\) 为在 \(\varphi\) 之下 \(\mathbf{B}_n\) 的一条直径的象；这个象是一条射影直线{]}。由此推出，若 \(n > \mathbf{1}\) ，则 \(\mathbf{P}_n\) 既不同胚于 \(\mathbf{S}_n\) ，也不同胚于 \(\mathbf{B}_n\) （利用 § 3 的习题 3 与 § 2 的习题 3）。
\item
  若把 \(\mathbf{R}^2\) 看作嵌在 \(\tilde{\mathbf{R}} \times \tilde{\mathbf{R}}\) 中，而把 \(\mathbf{R}\) 看作嵌在 \(\tilde{\mathbf{R}}\) 中，证明映射 \((x, y) \to x + y\) （ \(\mathbf{R}^2\) 到 \(\mathbf{R}\) 中的）可以连续延拓到点 \((\infty, a)\) 与 \((a, \infty)\) （它们属于 \(\tilde{\mathbf{R}} \times \tilde{\mathbf{R}}\) ， \(a\) 取一切有限值），但不能连续延拓到 \((\infty, \infty)\) 。若把 \(\mathbf{R}^2\) 看作嵌在 \(\mathbf{P}_2\) 中（ \(\mathbf{R}\) 仍看作嵌在 \(\tilde{\mathbf{R}}\) 中），则 \(x + y\) 可以连续延拓到无穷远直线上除齐次坐标为 \((1, -1, 0)\) 的点以外的所有点。
\end{enumerate}

对乘积 xy 陈述并证明这些结果的类似命题。

\begin{enumerate}
\def\labelenumi{\arabic{enumi})}
\setcounter{enumi}{5}
\item
  若 \(n \geqslant 0\) ，设 \(\mathfrak{F}\) 是 \(\mathbf{P}_n\) 的所有闭子集所成的集合，并按第 II 章 § 2 习题 6 b) 的程序赋予由 \(\mathbf{P}_n\) 的一致结构诱导的一致结构。证明：集合 \(\mathbf{P}_{n,p}\) （由 \(p \geqslant 0\) 维射影线性簇组成）在 \(\mathfrak{F}\) 中是闭的，且在 \(\mathbf{P}_{n,p}\) 上由 \(\mathfrak{F}\) 的拓扑诱导的拓扑与第 5 小节定义 2 中所定义的拓扑相同。{[}要定义 \(\mathbf{P}_n\) 的一致结构的围域，可把 \(\mathbf{P}_n\) 看作 \(\mathbf{S}_n\) 的商空间（第 1 小节命题 2），并取 \(\mathbf{S}_n\) 的被半径趋于 o 的球所作的有限覆盖。{]}
\item
  \begin{enumerate}
  \def\labelenumii{\alph{enumii})}
  \tightlist
  \item
    证明：在第 5 小节的记号下，空间 \(\mathbf{L}_{n,p}\) 同胚于 \(\mathbf{V}_{n,p}\) 与 \(\mathbb{R}^{p(p+1)/2}\) 的乘积。注意每个矩阵 \(X\in\mathbf{L}_{n,p}\) 都可以唯一地表示成 \(U\) . \(\Upsilon\) 的形式，其中 \(\Upsilon\in\mathbf{V}_{n,p}\) ，而 \(U=(u_{ij})\) 满足 \(u_{ij}=0\) （对 \(i<j\) ），且 \(u_{ii}>0\) （对 \(\mathbf{i}\leqslant i\leqslant p\) ）。令 \(\Upsilon=f(X)\) 。
  \end{enumerate}
\end{enumerate}

\begin{enumerate}
\def\labelenumi{\alph{enumi})}
\setcounter{enumi}{1}
\tightlist
\item
  在第 5 小节命题 8 的记号下，证明 \(f\) 是 \(\mathbf{B}_{\sigma}\) 到 \(f(\mathbf{B}_{\sigma})\) 上的一个同胚{[}注意 \(X \to f(X_{\sigma}^{-1}X)\) 是 \(\mathbf{A}_{\sigma}\) 到 \(f(\mathbf{B}_{\sigma})\) 上的连续映射，且 \(f(X_{\sigma}^{-1}X)\) 与 \(X \mod \Delta_{n,p}\) 属于同一类；然后应用引理 2{]}。由此推出， \(\mathbf{D}_{\sigma}\) （即 \(\mathbf{A}_{\sigma}\) 与 \(\mathbf{V}_{n,p}\) 的交）同胚于 \(\mathbf{V}_{p,p}\) 与 \(\mathbb{R}^{p(n - p)}\) 的乘积{[}属于 \(\mathbf{D}_{\sigma}\) 的每个矩阵都可以唯一地表示成 \(\mathbf{V}_{p,p}\) 中一个矩阵与 \(f(\mathbf{B}_{\sigma})\) 中一个矩阵的乘积{]}。
\end{enumerate}

\begin{enumerate}
\def\labelenumi{\arabic{enumi})}
\setcounter{enumi}{7}
\tightlist
\item
  设 \(g\) 是这样的映射：对每个矩阵 \(X \in \mathbf{L}_{n,p}\) ， \(X'g(X)\) 是由前 \(q\) 行（取自 \(X\) ）构成的矩阵（ \(q < p\) ）。证明： \(g\) 在 \(\mathbf{V}_{n,p}\) 上的限制是 \(\mathbf{V}_{n,p}\) 到 \(\mathbf{V}_{n,q}\) 上的连续开映射{[}利用习题 7 a)，注意若 \(X = U.Y\) 则
\end{enumerate}

\[
g (X) = U ^ {\prime}. g (\Upsilon),
\]

其中 \(U'\) 是由 U 去掉指标 \textgreater q 的行与列所得的矩阵；另一方面，注意 f 是开映射{]}。

\begin{enumerate}
\def\labelenumi{\arabic{enumi})}
\setcounter{enumi}{8}
\item
  证明： \(\mathbf{V}_{n,p}\) 在 \(p < n\) 时是连通的（利用习题 8 与第 I 章 § 11 第 3 小节命题 7）。
\item
  在射影空间 \(P_{n}\) 中，设 \(H_{n,p}\) 是由方程
\end{enumerate}

\[
x _ {0} ^ {2} + x _ {1} ^ {2} + \dots + x _ {p - 1} ^ {2} - x _ {p} ^ {2} - \dots - x _ {n} ^ {2} = 0 \quad (\mathrm{I} \leqslant p \leqslant n).
\]

所定义的“二次超曲面”。证明： \(H_{n,1}\) 与 \(H_{n,n}\) 都同胚于 \(S_{n-1}\) 。若 \(2 \leqslant p \leqslant n - 1\) ，则 \(H_{n,p}\) 同胚于把乘积 \(S_{p-1} \times S_{n-p}\) （视为 \(R^{p} \times R^{n-p+1}\) 的子空间）的每个点与其对径点等同起来所得到的空间。 \(H_{n,p}\) 的每个点都有一个同胚于 \(R^{n-1}\) 的开邻域。

证明： \(H_{3,2}\) 同胚于 \(S_{1} \times S_{1}\) {[}把 \(S_{1} \times S_{1}\) 与 \(T \times T\) 等同起来， \(S_{1} \times S_{1}\) 的一对对径点与点 \((u, v)\) 和 \((1/2 + u, 1/2 + v)\) （属于 \(T \times T\) ）相等同；然后考虑映射 \((u, v) \to (u + v, u - v)\) （ \(T \times T\) 到其自身的映射）{]}。

\begin{enumerate}
\def\labelenumi{\arabic{enumi})}
\setcounter{enumi}{10}
\tightlist
\item
  在射影空间 \(P_{n}\) 中，设 \(C_{n,p}\) 是由方程
\end{enumerate}

\[
x _ {1} ^ {2} + x _ {2} ^ {2} + \dots + x _ {p} ^ {2} - x _ {p + 1} ^ {2} - \dots - x _ {n} ^ {2} = 0 \quad (\mathrm{I} \leqslant p \leqslant n - \mathrm{I}).
\]

证明： \(\{0\}\) 在 \(C_{n,p}\) 中的补集同胚于 \(\mathbf{R} \times H_{n-1,p}\) （沿用习题 10 的记号）。

